Increasing query width during prompt processing or speculative verification can change a quantized language model's numerical behavior. We study how to reorganize this work while preserving a fixed deployment's outputs and state. Native operand-replay experiments show that a causal mask does not determine attention arithmetic: extending a physical KV range can change the result even when the added probabilities are exactly zero. Query width also changes activation quantization, reduction organization and device placement, while producer history and output ownership connect these effects across inference phases. We use these findings to construct reference-derived execution plans. The runtime derives compatible query groups from live cache geometry, places cuts around whole graphs or within attention, and exposes independent arithmetic groups to the GPU while retaining each group's reference computation. One resident owner preserves requested outputs and the committed frontier. CUDA and Vulkan implementations cover Qwen, Llama and SmolLM2, including mixed Q4/Q6 quantization at 32B. Compatible speculative applications reduce Linux/A5000 resident service time by 11.52\%--25.54\% against the fastest measured fixed target-only controls. Balanced experiments attribute 0.80\%--1.64\% to Q8 projection packing; the smaller contribution of local cuts follows the available cache-boundary geometry. Complete-request measurements identify when parallel execution improves service.
